% ----------------------
\subsection{Section 120 (8 July 1838)}
% ----------------------
	\label{DC120}
	
	% -----
	\subsubsection{Historical Background}
	% -----
		\label{DC120-HB}
		
		% --
		\paragraph{JSP}
		% --
		
			``On Sunday, 8 July 1838, JS dictated five revelations, each of which concerned church leadership or finance; one of these revelations regarded the disposition of property donated to the church. The management of such property had evolved over the years. Initially, the disposition of property was within the purview of the church's bishops. Then, from 1832 to 1834 the United Firm managed church assets. In the mid-1830s, after presidencies and high councils were established in Ohio and Missouri, they also became involved in overseeing the use of church assets. In 1837 William W. Phelps and John Whitmer, counselors in the Zion presidency, began making significant financial decisions without the input of the high council or the bishopric. High council members in particular resented the exclusion, and council members compelled Phelps and Whitmer to include the council and bishopric in future financial decisions.
			
			``The 8 July 1838 revelation regarding the management of donated property was the fourth of the five revelations from that day that George W. Robinson copied into JS's journal. JS apparently dictated these revelations in a leadership meeting held in the morning before the day's worship services, and at some point in July, Robinson copied the revelations into JS's journal as part of the entry for 8 July. This entry, which consists almost entirely of revelation transcripts, appears in a gap in regular journal keeping. Robinson apparently did not resume making regular journal entries until late July, suggesting that he may not have copied the revelations into the journal before then.
			
			``The third revelation directed the Latter-day Saints to consecrate surplus property to the church as `the beginning of the tithing of my people' and then to donate one-tenth of their interest annually. Robinson's headnote for the fourth revelation, featured here, states that it was given in reference to `the disposition of the properties tithed, as named in the preceeding revelation.' This short revelation directed that the donated property be managed by a council consisting of the First Presidency, Bishop Edward Partridge and his counselors in the Zion bishopric, and the Zion high council---all acting together under the inspiration of God. Whereas the first three 8 July revelations included in the journal entry for that day were read to the congregation of Saints at the worship service later in the day, it is unclear whether this fourth revelation was also read at that time. On 26 July 1838, the council that was called for in this revelation met as directed.''
					
	% -----
	\subsubsection{Versions}
	% -----
		\label{DC120-V}
		
		See the \href{https://drive.google.com/file/d/1goAvNz8jS4R8slYfMG_db55Ty2z_vmgK/view?usp=sharing}{sections comparison} document.
	
		\begin{compactdesc}
			\item[JSP] \hyperref[EMS-1838-JS-8-JUL-1838-D]{8 July 1838}
			\item[BC] ---
			\item[DC 1835] ---
			\item[DC 1844] ---
			\item[TS] ---
			\item[DN] 3:10 (2 April 1853), 37
			\item[MS] 16:12 (25 March 1854), 183
		\end{compactdesc}

	% -----
	\subsubsection{Section 120:1} % *DC120-1 
	% -----
		\label{DC120-1}

		VERILY, thus saith the Lord, the time is now come, that it shall be disposed of by a council, composed of the First Presidency of my Church, and of the bishop and his council, and by my high council; and by mine own voice unto them, saith the Lord.  Even so.  Amen.

		% \begin{framed}
		% \scriptsize
			% \begin{compactdesc}
				% \item[JSP] 
				% % \item[BC] 
				% % \item[]
			% \end{compactdesc}
		% \end{framed}